\assignment{1}{Mon 29 Aug 2022 16:23}{assignment}

\begin{enumerate}
\item Observer A, who is at rest in the laboratory, is studying a particle that is moving through the laboratory at a speed of $0.624 c$ and determines its lifetime to be \SI{159}{\nano\second}.
	\begin{enumerate}
		\item Observer A places markers in the laboratory at the locations where the particle is produced and where it decays. How far apart are those markers in the laboratory?
\begin{proof}[Solution]
	\begin{align*}
		x_A = v_A t_A &= \SI{0.624}{c} \times \SI{159}{ns}\\
		&= 0.624\times \SI{3e8}{m\per s} \times \SI{1.59e-7}{s} \\
		&= \SI{29.7648}{m}
	\end{align*}
\end{proof}

	\item Observer B, who is traveling parallel to the particle at a speed of $0.624 c$, observes the particle to be at rest and measures its lifetime to be \SI{124}{\nano\second}. According to B, how far apart are the two markers in the laboratory?
\begin{proof}[Solution]
	\begin{align*}
		x_B = v_B t_B &= \SI{0.624}{c} \times \SI{124}{ns} \\
		&= \SI{1.872e8}{m \per s}\times  \SI{1.24e-7}{s} \\
		&= \SI{23.21}{m}
	.\end{align*}
\end{proof}

	\end{enumerate}
\item For a gas at thermal equilibrium at absolute temperature $\mathrm{T}$, the distribution of molecular energies is given by the Maxwell-Boltzmann distribution
\[
N(E)=\frac{2 N}{\sqrt{\pi}} \frac{1}{(k T)^{3 / 2}} E^{1 / 2} e^{-E / k T}
\]
where $N$ is the total number of molecules and $N(E)$ is the distribution function defined so that $N(E) d E$ is the number of molecules in the energy interval between $E$ and $E+d E$. Show that the most probably molecular energy is $\frac{1}{2} k T$.

\begin{proof}[Solution]
	We assume that the distribution is continuous and unimodal. It suffices to find an extremum value for $N(E)$.
 From calculus we know that a continous function reaches extremum iff its first derivative is zero at that point.\begin{align*}
	N'(E)&= 0 \\
	\frac{\partial }{\partial E} \left( \frac{2 N}{\sqrt{\pi}} \frac{1}{(k T)^{3 / 2}} E^{1 / 2} e^{-E / k T} \right) &= 0 \\
	\frac{\partial }{\partial E} \left( E^{1 / 2} e^{-E / k T} \right) &= 0 \\
	\frac{1}{2}E^{-\frac{1}{2}}e^{-\frac{E}{kT}} + E^{\frac{1}{2}}\left( -\frac{1}{KT} \right) e^{-\frac{E}{kT}} &= 0 \\
	\frac{1}{2}E^{-\frac{1}{2}} + E^{\frac{1}{2}}\left( -\frac{1}{KT} \right)  &= 0 \\
	\frac{1}{2} + E\left( -\frac{1}{KT} \right)  &= 0 \\
	E &= \frac{\frac{1}{2}}{\frac{1}{kT}} \\
	E &= \frac{1}{2}kT 
.\end{align*}
\end{proof}


\item For a molecule of Oxygen $\left(\mathrm{O}_{2}\right)$ at room temperature \SI{300}{\kelvin} aligned on the z-axis, calculate the average angular velocity for rotations about the $\mathrm{x}$-axis. The distance between the O atoms in the molecule is \SI{0.121}{\nano\metre}.
\begin{proof}[Solution]
	We assume that classical statistical mechanics work for the problem. The rotation about $x$-axis corresponds to one degree of freedom. By theorem of equipartition of energy, it takes $\frac{1}{2}k_B T$ of energy. Therefore,
\begin{align*}
	E_{\text{rot}}&= \frac{1}{2}k_B T \\
	&= \frac{1}{2} \SI{1.38e-23}{m\squared \kilogram \per s \squared \kelvin} \times  \SI{300}{K} \\
	&=  \SI{2.07e-21}{J}
.\end{align*}
We also assume that we can view an oxygen molecule as two oxygen atoms connected by a lightweight and rigid rod. We may consider this rod to be rigid since we are only interested in rotational motion, not vibrational.

Use the known data that an oxygen atom has mass  \SI{15.999}{\atomicmassunit}. Its moment of inertia around $x$ axis at its center of mass is then
\begin{align}
I = \frac{1}{2} m x^2 &= \frac{1}{2} \SI{15.999}{\atomicmassunit} (\SI{0.121}{\nano\metre})^2 \\
&=  \SI{1.9442e-46}{\kilogram \square \metre}
.\end{align}
Finally we solve the equation for rotational kinetic energy:

\begin{align*}
	E_{\text{rot}} &= \frac{1}{2} I \omega^2\\
	\omega^2 &= \SI{2.1294e+25}{}\\
	\omega &= \SI{4.6146e+12}{\radian \per \second}\\
	\omega &= \SI{7.3442e+11}{rounds \per \second}
.\end{align*}

They are turning so damn fast!
\end{proof}

\end{enumerate}
